\section{AHSME 1989}
        \begin{enumerate}
        	\item (\textit{AHSME 1989}) $(-1)^{5^{2}}+1^{2^{5}}=$


 $\textrm{(A)}\ -7\qquad\textrm{(B)}\ -2\qquad\textrm{(C)}\ 0\qquad\textrm{(D)}\ 1\qquad\textrm{(E)}\ 57$


 
 \textbf{Answer: }C



\textbf{Solution}

$(-1)^{5^2} + 1^{2^5} = (-1)^{25}+1^{32}=-1+1=\boxed{\textbf{(C)}\ 0}$.

 


	\item (\textit{AHSME 1989}) $\sqrt{\frac{1}{9}+\frac{1}{16}}=$


 $\textrm{(A)}\ \frac{1}5\qquad\textrm{(B)}\ \frac{1}4\qquad\textrm{(C)}\ \frac{2}7\qquad\textrm{(D)}\ \frac{5}{12}\qquad\textrm{(E)}\ \frac{7}{12}$


 
 \textbf{Answer: }D



\textbf{Solution}

$\sqrt{\frac{1}{9}+\frac{1}{16}}=\sqrt{\frac{25}{144}}=\boxed{\textbf{(D)}\ \frac{5}{12}}$.

 


 


	\item (\textit{AHSME 1989}) A square is cut into three rectangles along two lines parallel to a side, as shown. If the perimeter of each of the three rectangles is $24$, then the area of the original square is


 
                    \begin{figure}[H]
                        centering
                        \includegraphics[width=0.5\textwidth]{Problem/AHSME_I_1989_3_Problem_Figure_1.png}
                    \end{figure}
                    


 $\textrm{(A)}\ 24\qquad\textrm{(B)}\ 36\qquad\textrm{(C)}\ 64\qquad\textrm{(D)}\ 81\qquad\textrm{(E)}\ 96$


 
 \textbf{Answer: }D



\textbf{Solution}

Let $x$ be the width of one of the three rectangles.  Then the square's side length is $3x$, so the perimeter of each rectangle is $2(3x + x) = 8x$. Since $8x=24$, $x=3$ and the square's side length is $3 \cdot 3 = 9$. Therefore, square's area is $9^2 = \boxed{\textbf{(D)}\ 81}$.

 


	\item (\textit{AHSME 1989}) In the figure, $ABCD$ is an isosceles trapezoid with side lengths $AD=BC=5$, $AB=4$, and $DC=10$. The point $C$ is on $\overline{DF}$ and $B$ is the midpoint of hypotenuse $\overline{DE}$ in right triangle $DEF$. Then $CF=$


 
                    \begin{figure}[H]
                        centering
                        \includegraphics[width=0.5\textwidth]{Problem/AHSME_I_1989_4_Problem_Figure_1.png}
                    \end{figure}
                    


 $\textrm{(A)}\ 3.25\qquad\textrm{(B)}\ 3.5\qquad\textrm{(C)}\ 3.75\qquad\textrm{(D)}\ 4.0\qquad\textrm{(E)}\ 4.25$


 
 \textbf{Answer: }D



\textbf{Solution 1}


                    \begin{figure}[H]
                        centering
                        \includegraphics[width=0.5\textwidth]{Solution/AHSME_I_1989_4_Solution_Figure_1.png}
                    \end{figure}
                    

 Let $X$ and $Y$ be the feet of the perpendiculars from $A$ and $B$ to $\overline{DF}$, respectively. Then $\triangle DEF \sim \triangle DBY$, and the sides of $\triangle DEF$ are twice as long as those of $\triangle DBY$.

 So $XY = AB = 4$, and $DX = YC = \frac{10-4}{2} = 3$.  Thus $DY = 4 + 3 = 7$ and $DF = 2 \cdot DY = 14$.  Finally, $CF = 14 - 10 = \boxed{\textbf{(D)}\ 4.0}$.

 



\textbf{Solution 2}


                    \begin{figure}[H]
                        centering
                        \includegraphics[width=0.5\textwidth]{Solution/AHSME_I_1989_4_Solution_Figure_2.png}
                    \end{figure}
                    

 Draw $\overline{BF}$.  Since $B$ is the midpoint of $\overline{DE}$, $BD = BF = BE$.  So $\triangle BDF$ is isosceles, and $\angle BFD = \angle BDF = \angle ABD$.  Also, $\angle A = \angle ABC = \angle BCF$.  Therefore, $\triangle DAB \cong \triangle BCF$, so $CF = AB = \boxed{\textbf{(D)}\ 4.0}$.

 


	\item (\textit{AHSME 1989}) Toothpicks of equal length are used to build a rectangular grid as shown. If the grid is 20 toothpicks high and 10 toothpicks wide, then the number of toothpicks used is


 
                    \begin{figure}[H]
                        centering
                        \includegraphics[width=0.5\textwidth]{Problem/AHSME_I_1989_5_Problem_Figure_1.png}
                    \end{figure}
                    


 $\textrm{(A)}\ 30\qquad\textrm{(B)}\ 200\qquad\textrm{(C)}\ 410\qquad\textrm{(D)}\ 420\qquad\textrm{(E)}\ 430$


 
 \textbf{Answer: }E



\textbf{Solution 1}

The grid consists of $20+1=21$ horizontal $10$-toothpick segments, and $10+1=11$ vertical $20$-toothpick segments, for a total of $21\cdot10+11\cdot20=\boxed{\textbf{(E)}\ 430}$ toothpicks.

 



\textbf{Solution 2}


                    \begin{figure}[H]
                        centering
                        \includegraphics[width=0.5\textwidth]{Solution/AHSME_I_1989_5_Solution_Figure_1.png}
                    \end{figure}
                    

 Divide the figure into $20 \cdot 10 = 200$ L-shaped joints, the right edge, and the bottom edge.  It becomes clear that there are $2\cdot200+20+10=\boxed{\textbf{(E)}\ 430}$ toothpicks.

 


	\item (\textit{AHSME 1989}) If $a,\,b>0$ and the triangle in the first quadrant bounded by the coordinate axes and the graph of $ax+by=6$ has area $6$, then $ab=$


 $\textrm{(A)}\ 3\qquad\textrm{(B)}\ 6\qquad\textrm{(C)}\ 12\qquad\textrm{(D)}\ 108\qquad\textrm{(E)}\ 432$


 
 \textbf{Answer: }A



\textbf{Solution}

Set $y=0$ to find that the $x$-intercept of the line is $\frac{6}{a}$. Set $x=0$ to find that the $y$-intercept of the line is $\frac{6}{b}$. 

 In terms of $a$ and $b$, the triangle's area is $\frac{1}{2}\cdot\frac{6}{a}\cdot\frac{6}{b} = \frac{18}{ab}$. So $\frac{18}{ab} = 6$, and thus $ab=\boxed{\textbf{(A)}\ 3}$.

 -$\LaTeX$ by Kevinliu08, edited by j314andrews

 



\textbf{Solution 2}

We find the $y$-intercept of the line to be $6$, which means the two legs of the triangle are $6$ and $2$, respectively. We then conclude that the line passes through the points $(0,6)$ and $(2,0)$; when we plug these points into the equation, we get that $a=3$ and $b=1$. We multiply these together and see that the answer is $\boxed{\textbf{(A)}\ 3}$.

 


	\item (\textit{AHSME 1989}) In $\triangle ABC$, $\angle A = 100^\circ$, $\angle B = 50^\circ$, $\angle C = 30^\circ$, $\overline{AH}$ is an altitude, and $\overline{BM}$ is a median. Then $\angle MHC=$


 
                    \begin{figure}[H]
                        centering
                        \includegraphics[width=0.5\textwidth]{Problem/AHSME_I_1989_7_Problem_Figure_1.png}
                    \end{figure}
                    


 $\textrm{(A)}\ 15^\circ\qquad\textrm{(B)}\ 22.5^\circ\qquad\textrm{(C)}\ 30^\circ\qquad\textrm{(D)}\ 40^\circ\qquad\textrm{(E)}\ 45^\circ$


 
 \textbf{Answer: }C



\textbf{Solution 1}


                    \begin{figure}[H]
                        centering
                        \includegraphics[width=0.5\textwidth]{Solution/AHSME_I_1989_7_Solution_Figure_1.png}
                    \end{figure}
                    Since $\overline{BM}$ is a median, $AM = MC$. Let $N$ be the foot of the perpendicular from $M$ to $\overline{HC}$. Since $MN \parallel AH$, $HN = NC$, and thus $\triangle MNH \cong \triangle MNC$. So $\angle MHC=\angle C=\boxed{\textbf{(C)}\ 30^\circ}$. 

 



\textbf{Solution 2}

Since $M$ is the midpoint of hypotenuse $\overline{AC}$ of $\triangle AHC$, $M$ is the circumcenter of $\triangle AHC$ and thus $MA = MH = MC$. So $\triangle MHC$ is isosceles and $\angle MHC = \angle C = \boxed{\textbf{(C)}\ 30^\circ}$.

 


	\item (\textit{AHSME 1989}) For how many integers $n$ between $1$ and $100$ does $x^{2}+x-n$ factor into the product of two linear factors with integer coefficients?


 $\textrm{(A)}\ 0\qquad\textrm{(B)}\ 1\qquad\textrm{(C)}\ 2\qquad\textrm{(D)}\ 9\qquad\textrm{(E)}\ 10$


 
 \textbf{Answer: }D



\textbf{Solution 1}

For $x^2+x-n$ to factor into a product of two linear factors, we must have $x^2+x-n = (x + a)(x + b)$, where $a$ and $b$ are integers.

 By expansion of the product of the linear factors and comparison to the original quadratic,

 $ab = -n$

 $a + b = 1$.

 The only way for this to work if $n$ is a positive integer is if $a = -b +1$.

 Here are the possible pairs:

 $a = -1, b = 2$

 $a = -2, b = 3$

 $\vdots$

 $a = -9, b = 10$

 

 This gives us 9 integers for $n$, $\boxed{\text{D}}$.

 



\textbf{Solution 2}

By the quadratic formula, the roots of $x^2+x-n$ are $\frac{-1 \pm \sqrt{4n+1}}{2}$.  For $x^2+x-n$ to factor into two linear factors with integer coefficients, both roots must be integers and thus $4n+1$ must be an odd perfect square.

 All odd perfect squares have remainder $1$ when divided by $4$, so $4n+1$ may be any odd perfect square between $4\cdot 1 + 1 = 5$ and $4 \cdot 100 + 1 = 401$.  Therefore, $4n+1$ could be one of $3^2, 5^2, ..., 17^2, 19^2$, which gives $\boxed{\textbf{(D)}\ 9}$ values for $n$.

 



\textbf{Solution 3}

Note that $x^2+x-n$ factors into two linear factors with integer coefficients if and only if both its roots are integers.  Also, since the product of the roots is $-n$, one root is positive while the other is negative.

 Let $a$ be the positive root.  Then $n = a^2+a = a(a+1)$.  Since $n$ must be between $1$ and $100$, $a$ may range from $1$ to $9$, giving a total of $\boxed{\textbf{(D)}\ 9}$ possible values for $n$.

 -j314andrews

 


	\item (\textit{AHSME 1989}) Mr. and Mrs. Zeta want to name baby Zeta so that its monogram (first, middle, and last initials) will be in alphabetical order with no letter repeated. How many such monograms are possible? 


 $\textrm{(A)}\ 276\qquad\textrm{(B)}\ 300\qquad\textrm{(C)}\ 552\qquad\textrm{(D)}\ 600\qquad\textrm{(E)}\ 15600$


 
 \textbf{Answer: }B



\textbf{Solution}

We see that for any combination of two distinct letters other than Z (as the last name will automatically be Z), there is only one possible way to arrange them in alphabetical order, thus the answer is just $\dbinom{25}{2}=\boxed{\mathrm{(B)\,300}}$.

 


	\item (\textit{AHSME 1989}) Consider the sequence defined recursively by $u_{1}= a$ (any positive integer), and $u_{n+1}=\frac{-1}{u_{n}+1}$,  $n = 1,2,3,\cdots$. For which of the following values of $n$ must $u_{n}=a$?


 $\textrm{(A)}\ 14\qquad\textrm{(B)}\ 15\qquad\textrm{(C)}\ 16\qquad\textrm{(D)}\ 17\qquad\textrm{(E)}\ 18$


 
 \textbf{Answer: }C



\textbf{Solution}

Repeatedly applying the function, and simplifying, we get 
 \[a,\quad-\frac1{a+1},\quad-\frac{a+1}a,\]and then $a$ again. So $a$ must appear at every third term after $u_1$. The only option given of the form $1+3k$ is $\boxed{\mathrm{(C)}\,16}$.

 


	\item (\textit{AHSME 1989}) Let $a$, $b$, $c$, and $d$ be positive integers with $a < 2b$, $b < 3c$, and $c<4d$. If $d<100$, the largest possible value for $a$ is


 $\textrm{(A)}\ 2367\qquad\textrm{(B)}\ 2375\qquad\textrm{(C)}\ 2391\qquad\textrm{(D)}\ 2399\qquad\textrm{(E)}\ 2400$


 
 \textbf{Answer: }A



\textbf{Solution}

Each of these integers is bounded above by the next one.

 
                \begin{itemize}
                    \item $d<100$, so the maximum $d$ is $99$.\item $c<4d\le396$, so the maximum $c$ is $395$.\item $b<3c\le1185$, so the maximum $b$ is $1184$.\item $a<2b\le2368$, so the maximum $a$ is $\boxed{2367}$.
                \end{itemize}


                Note that the statement $a<2b<6c<24d<2400$ is true, but does not specify the distances between each pair of values.

 


	\item (\textit{AHSME 1989}) The traffic on a certain east-west highway moves at a constant speed of 60 miles per hour in both directions. An eastbound driver passes 20 west-bound vehicles in a five-minute interval. Assume vehicles in the westbound lane are equally spaced. Which of the following is closest to the number of westbound vehicles present in a 100-mile section of highway?


 $\textrm{(A)}\ 100\qquad\textrm{(B)}\ 120\qquad\textrm{(C)}\ 200\qquad\textrm{(D)}\ 240\qquad\textrm{(E)}\ 400$


 
 \textbf{Answer: }C



\textbf{Solution 1}

At the beginning of the five-minute interval, say the eastbound driver is at the point $x=0$, and at the end of the interval at $x=5$, having travelled five miles. Because both lanes are travelling at the same speed, the last westbound car to be passed by the eastbound driver was just west of the position $x=10$ at the start of the five minutes. The first westbound car to be passed was just east of $x=0$ at that time. Therefore, the eastbound driver passed all of the cars initially in the stretch of road between $x=0$ and $x=10$. That makes $20$ cars in ten miles, so we estimate $200$ cars in a hundred miles.

 
                    \begin{figure}[H]
                        centering
                        \includegraphics[width=0.5\textwidth]{Solution/AHSME_I_1989_12_Solution_Figure_1.png}
                    \end{figure}
                    

 



\textbf{Solution 2}

Since the westbound vehicles are equally spaced and move at a constant speed, the eastbound driver passes a westbound vehicle approximately once every $\frac{5}{20} = \frac{1}{4}$ minute, or every $\frac{1}{240}$ hour.  

 Let $t = 0$ hours be the time eastbound driver passes its first westbound vehicle.  Then the eastbound driver passes its second westbound vehicle at approximately $t = \frac{1}{240}$ hour.  In this time, the eastbound and second westbound vehicles each traveled about $60 \cdot \frac{1}{240} = \frac{1}{4}$ miles, so they were about $\frac{1}{2}$ mile apart at $t = 0$. Since the first westbound and eastbound vehicles were at the same position at $t = 0$, the first and second westbound vehicles were also about $\frac {1}{2}$ mile apart.  

 Therefore, each pair of consecutive westbound vehicles are about $\frac{1}{2}$ mile apart, zo a $100$-mile section must contain about $\frac{100}{\frac{1}{2}} = \boxed{(\textbf{C})\ 200}$ of them.  

 -j314andrews

 


	\item (\textit{AHSME 1989}) Two strips of width 1 overlap at an angle of $\alpha$ as shown. The area of the overlap (shown shaded) is


 
                    \begin{figure}[H]
                        centering
                        \includegraphics[width=0.5\textwidth]{Problem/AHSME_I_1989_13_Problem_Figure_1.png}
                    \end{figure}
                    


 $\textrm{(A)}\ \sin\alpha\qquad\textrm{(B)}\ \frac{1}{\sin\alpha}\qquad\textrm{(C)}\ \frac{1}{1-\cos\alpha}\qquad\textrm{(D)}\ \frac{1}{\sin^{2}\alpha}\qquad\textrm{(E)}\ \frac{1}{(1-\cos\alpha)^{2}}$


 
 \textbf{Answer: }B



\textbf{Solution}


                    \begin{figure}[H]
                        centering
                        \includegraphics[width=0.5\textwidth]{Solution/AHSME_I_1989_13_Solution_Figure_1.png}
                    \end{figure}
                    

 Let $s$ be the side length of the rhombus.  Then by definition $\sin \alpha = \frac{1}{s}$, so $s = \frac{1}{\sin \alpha}$.

 Since the rhombus has height $1$, its area is $\boxed{\textbf{(B)}\ \frac{1}{\sin\alpha}}$

 


	\item (\textit{AHSME 1989}) $\cot 10+\tan 5 =$


 $\textrm{(A)}\ \csc 5\qquad\textrm{(B)}\ \csc 10\qquad\textrm{(C)}\ \sec 5\qquad\textrm{(D)}\ \sec 10\qquad\textrm{(E)}\ \sin 15$


 
 \textbf{Answer: }B



\textbf{Solution 1}

$\cot 10 +\tan 5=\frac{\cos 10}{\sin 10}+\frac{\sin 5}{\cos 5}=\frac{\cos10\cos5+\sin10\sin5}{\sin10\cos 5}=\frac{\cos(10-5)}{\sin10\cos5}=\frac{\cos5}{\sin10\cos5}=\boxed{\textbf{(B)}\csc 10}$.

 



\textbf{Solution 2 (Half Angle Formulas)}

$\cot 10 + \tan 5 = \frac{\cos 10}{\sin 10} + \frac{1 - \cos 10}{\sin 10} = \frac{1}{\sin 10} = \boxed{(\textbf{B})\csc 10}$.

 -j4andrews

 



\textbf{Solution 3 (Double Angle Formulas)}

$\cot 10 + \tan 5 = \frac{1}{\tan 10} + \tan 5 = \frac{1}{\frac{2 \tan 5}{1 - \tan^2 5}} + \tan 5 = \frac{1 - \tan^2 5}{2 \tan 5} + \tan 5 = \frac{1 - \tan^2 5}{2 \tan 5} + \frac{2 \tan^2 5}{2 \tan 5}$

 $= \frac{1 + \tan^2 5}{2 \tan 5} = \frac{\sec^2 5}{2 \tan 5} = \frac{1}{2} \sec^2 5 \cot 5 = \frac{1}{2} \cdot \frac{1}{\cos^2 5} \cdot \frac {\cos 5}{\sin 5} = \frac{1}{2 \sin 5 \cos 5} = \frac{1}{\sin 10} = \boxed{(\textbf{B})\csc 10}$.

 -j314andrews

 


	\item (\textit{AHSME 1989}) In $\triangle ABC$, $AB=5$, $BC=7$, $AC=9$, and $D$ is on $\overline{AC}$ with $BD=5$. Find the ratio of $AD:DC$. 


 
                    \begin{figure}[H]
                        centering
                        \includegraphics[width=0.5\textwidth]{Problem/AHSME_I_1989_15_Problem_Figure_1.png}
                    \end{figure}
                    


 $\textrm{(A)}\ 4:3\qquad\textrm{(B)}\ 7:5\qquad\textrm{(C)}\ 11:6\qquad\textrm{(D)}\ 13:5\qquad\textrm{(E)}\ 19:8$


 
 \textbf{Answer: }E



\textbf{Solution}


                    \begin{figure}[H]
                        centering
                        \includegraphics[width=0.5\textwidth]{Solution/AHSME_I_1989_15_Solution_Figure_1.png}
                    \end{figure}
                    Drop the altitude $h$ from $B$ through $AD$, and let $AD$ be $2x$. Then by Pythagoras 
 \[\begin{cases}h^2+x^2=25\\h^2+(9-x)^2=49\end{cases}\] and after subtracting the first equation from the second, $x=3\frac16$. Therefore the desired ratio is 
 \[\frac{6\frac13}{2\frac23}=\boxed{\frac{19}8}\]

 


 



\textbf{Solution 2 (Trig)}

Using laws of cosines on $\bigtriangleup ABC$ yields $49=25+81-2 \cdot 5 \cdot 9 \cdot \cos A \implies \cos A =\frac{19}{30}.$ Let $AD=x.$ Using laws of cosines on $\bigtriangleup ABD$ yields $25+x^2-2 \cdot 5 \cdot x \cdot \cos A = 25.$ Fortunately, we know that $\cos A =\frac{19}{30}.$ Plugging this information back into our equation yields $x=\frac{19}{3}.$ Then, we know that $DC=9-\frac{19}{3}=\frac{8}{3} \implies\frac{AD}{DC}=\boxed{\frac{19}{8}}.$

 



\textbf{Solution 3 (Heron's Theorem)}

By Heron's Theorem, the area of $\triangle ABC$ is $\sqrt{\frac{21}{2} \cdot \frac{11}{2} \cdot \frac{7}{2} \cdot \frac{3}{2}} = \frac{21\sqrt{11}}{4}$.  

 Let $E$ be the foot of the perpendicular from $B$ to $\overline{AC}$, and let $h = BE$. Then $\frac{9h}{2} = \frac{21\sqrt{11}}{4}$, so $h = \frac{7\sqrt{11}}{6}$.  

 By the Pythagorean theorem, $AE = ED = \sqrt{5^2 - \left(\frac{7\sqrt{11}}{6}\right)^2} = \frac{19}{6}$. Thus $AD = \frac{19}{3}$, $DC = 9 - \frac{19}{3} = \frac{8}{3}$, and $AD : DC = \boxed{\textbf{(E)}\ 19 : 8}$.  

 



\textbf{Solution 4 (Stewart's Theorem)}

Let $x = AD$.

 By Stewart's Theorem, $(BD)^2(AC) + (AC)(AD)(DC) = (AB)^2(CD) + (BC)^2(AD)$.  

 Substituting yields $5^2(9) + 9x(9-x) = 5^2 (9-x) +7^2(x)$ which simplifies to $-9x^2 + 81x + 225 = 225 + 24x$.  Collecting all terms on one side yields $9x^2 - 57x = 0$, which factors as $3x(3x-19)=0$.  Since $x$ is positive, $x = \frac{19}{3}$.

 Therefore, $DC = 9 - \frac{19}{3} = \frac{8}{3}$ and $AD : DC = \boxed{\textbf{(E)}\ 19 : 8}$.

 


	\item (\textit{AHSME 1989}) A lattice point is a point in the plane with integer coordinates. How many lattice points are on the line segment whose endpoints are (3,17) and (48,281)? (Include both endpoints of the segment in your count.)


 $\textrm{(A)}\ 2\qquad\textrm{(B)}\ 4\qquad\textrm{(C)}\ 6\qquad\textrm{(D)}\ 16\qquad\textrm{(E)}\ 46$


 
 \textbf{Answer: }B



\textbf{Solution}

The difference in the $y$-coordinates is $281 - 17 = 264$, and the difference in the $x$-coordinates is $48 - 3 = 45$.The gcd of 264 and 45 is 3, so the line segment joining $(3,17)$ and $(48,281)$ has slope
 \[\frac{88}{15}.\]The points on the line have coordinates
 \[\left(3+t,17+\frac{88}{15}t\right).\]If $t$ is an integer, the $y$-coordinate of this point is an integer if and only if $t$ is a multiple of 15. The points where $t$ is a multiple of 15 on the segment $3\leq x\leq 48$ are $3$, $3+15$, $3+30$, and $3+45$. There are 4 lattice points on this line. Hence the answer is $\boxed{B}$.

 


	\item (\textit{AHSME 1989}) The perimeter of an equilateral triangle exceeds the perimeter of a square by $1989$ cm. The length of each side of the triangle exceeds the length of each side of the square by $d$ cm. The square has perimeter greater than 0. How many positive integers are NOT a possible value for $d$?


 $\textrm{(A)}\ 0\qquad\textrm{(B)}\ 9\qquad\textrm{(C)}\ 221\qquad\textrm{(D)}\ 663\qquad\textrm{(E)}\ \text{infinitely many}$


 
 \textbf{Answer: }D



\textbf{Solution}

Let $s$ be the side length of the square.  Then the triangle has side length $s + d$, and the difference of their perimeters is $3(s + d) - 4s = 3d - s = 1989$.  

 So $d = \frac{s + 1989}{3} = \frac{s}{3} + 663$.  Since $s > 0$, $d > 663$ and thus the first $663$ positive integers are not possible values of $d$.  Any integer $n \geq 664$ is a possible value of $d$, since $s = 3(n - 663)$ would suffice.  Therefore, $\boxed{\textbf{(D)}\ 663}$ positive integers are not possible values of $d$.

 


	\item (\textit{AHSME 1989}) The set of all numbers x for which $x+\sqrt{x^{2}+1}-\frac{1}{x+\sqrt{x^{2}+1}}$ is a rational number is the set of all:


 $\textrm{(A)}\ \text{ integers }x\qquad\textrm{(B)}\ \text{ rational }x\qquad\textrm{(C)}\ \text{ real }x\qquad\textrm{(D)}\ x\text{ for which }\sqrt{x^{2}+1}\text{ is rational}\qquad\textrm{(E)}\ x\text{ for which }x+\sqrt{x^{2}+1}\text{ is rational }$


 
 \textbf{Answer: }B



\textbf{Solution}

Rationalizing the denominator of $\frac{1}{x+\sqrt{x^2+1}}$, it simplifies: $\frac{1}{x+\sqrt{x^2+1}}$ = $\frac{x-\sqrt{x^2+1}}{x^2-(x^2+1)}$ = $\frac{x-\sqrt{x^2+1}}{-1}$ = $-(x-\sqrt{x^2+1})$. Substituting this into the original equation, we get $x + \sqrt{x^2+1} - (-(x-\sqrt{x^2+1})) = x+\sqrt{x^2+1} + x - \sqrt{x^2+1} = 2x$. $2x$ is only rational if $x$ is rational $\mathrm{(B)}$

 


	\item (\textit{AHSME 1989}) A triangle is inscribed in a circle. The vertices of the triangle divide the circle into three arcs of lengths $3$, $4$, and $5$. What is the area of the triangle?


 $\textrm{(A)}\ 6\qquad\textrm{(B)}\ \frac{18}{\pi^{2}}\qquad\textrm{(C)}\ \frac{9}{\pi^{2}}\left(\sqrt{3}-1\right)\qquad\textrm{(D)}\ \frac{9}{\pi^{2}}\left(\sqrt{3}+1\right)\qquad\textrm{(E)}\ \frac{9}{\pi^{2}}\left(\sqrt{3}+3\right)$


 
 \textbf{Answer: }E



\textbf{Solution}

The three arcs make up the entire circle, so the circumference of the circle is $3+4+5=12$ and the radius is $\frac{12}{2\pi}=\frac{6}{\pi}$. Also, the lengths of the arcs are proportional to their corresponding central angles. Thus, we can write the values of the arcs as $3\theta$, $4\theta$, and $5\theta$ for some $\theta$. By Circle Angle Sum, we obtain $3\theta+4\theta+5\theta=360$. Solving yields $\theta=30$. Thus, the angles of the triangle are $90$, $120$, and $150$. Using $[ABC]=\frac{1}{2}ab\sin{C}$, we obtain $\frac{r^2}{2}(\sin{90}+\sin{120}+\sin{150})$. Substituting $\frac{6}{\pi}$ for $r$ and evaluating yields $\frac{9}{\pi^2}(\sqrt{3}+3)\implies{\boxed{E}}$.

 
-Solution by \textbf{thecmd999}

 



\textbf{Solution 2 (no trig)}

Since we know that the triangle is inscribed, we can use the properties of inscribed angles to solve the problem. At first, we can ignore actual lengths in this problem, and use the $3$, $4$, and $5$ as ratios instead of the actual numbers. There are $360$ degrees in a circle, so $3x + 4x + 5x = 360$, where we find $x$ to be $30$ degrees. Thus, the arc measures in degrees becomes $90$, $120$, and $150$. With this knowledge, the angle that cuts off the arc of length $3$ is $45$ degrees (an inscribed angle is $\frac{1}{2}$ the measure of the arc that it intercepts). Similarly, the angle cutting off length $4$ is $60$ degrees, and the last angle of the triangle is $75$ degrees. 

 Upon observing this figure, we can draw an altitude in the triangle, breaking it into two triangles, one being a 30-60-90 and the other being a 45-45-90. For now, we can set the leg of the 30-60-90 angle that is adjacent to the 60 degree angle as "x." This makes the other leg of the 30-60-90 degree angle to be x√3, and since that happens to be the altitude and also part of the 45-45-90 triangle, it makes both legs of the 45-45-90 triangle to be x√3. We sum together the areas of both triangles and get (x\textasciicircum{}2√3 + 3x\textasciicircum{}2)/2, or x\textasciicircum{}2(√3+3)/2.

 Now it's time to actually find the value of x and plug it in. In order to find the value of x, we can use the length of the radius, which we identify as "r." Now, looking back at what we had earlier about the arcs, the arc of length 3 is 90 degrees, which is one-fourth of the circumference. The circumference itself is equal to 2rπ. So we can write this equation: 3 = 2rπ/4, making r = 6/π. Now we must find the relationship of r to x, but we don't have to look far to find it. In fact, if we draw the center of the circle and connect two radii from it to points of arc length 3, we see that it creates another 45-45-90 triangle, where r is the length of the leg, and as we see from earlier, the length of that hypotenuse is 2x. Since the hypotenuse is equal to the leg multiplied by √2 in such a triangle, we write this equation: 2x = 6√2/π, and x = 3√2/π. Now, we plug our new value of x into the equation x\textasciicircum{}2(√3+3)/2. Doing so, we get $\frac{9}{\pi^2}(\sqrt{3}+3)\implies{\boxed{E}}$.

 


	\item (\textit{AHSME 1989}) Let $x$ be a real number selected uniformly at random between 100 and 200. If $\lfloor {\sqrt{x}} \rfloor = 12$, find the probability that $\lfloor {\sqrt{100x}} \rfloor = 120$. ($\lfloor {v} \rfloor$ means the greatest integer less than or equal to $v$.)


 $\text{(A)} \ \frac{2}{25} \qquad \text{(B)} \ \frac{241}{2500} \qquad \text{(C)} \ \frac{1}{10} \qquad \text{(D)} \ \frac{96}{625} \qquad \text{(E)} \ 1$


 
 \textbf{Answer: }B



\textbf{Solution}

Since $\lfloor\sqrt{x}\rfloor=12$, $12\leq\sqrt{x}<13$ and thus $144\leq x<169$.

 The successful region is when $120\leq10\sqrt{x}<121$ in which case $12\leq\sqrt{x}<12.1$  Thus, the successful region is when
 \[144\leq x<146.41\]

 The successful region consists of a 2.41 long segment, while the total possibilities region is 25 wide.  Thus, the probability is
 \[\frac{2.41}{25}=\boxed{\frac{241}{2500}};\;\boxed{B}.\]

 


	\item (\textit{AHSME 1989}) A square flag has a red cross of uniform width with a blue square in the center on a white background as shown. (The cross is symmetric with respect to each of the diagonals of the square.) If the entire cross (both the red arms and the blue center) takes up 36\% of the area of the flag, what percent of the area of the flag is blue?


 
                    \begin{figure}[H]
                        centering
                        \includegraphics[width=0.5\textwidth]{Problem/AHSME_I_1989_21_Problem_Figure_1.png}
                    \end{figure}
                    


 $\text{(A)}\ .5\qquad\text{(B)}\ 1\qquad\text{(C)}\ 2\qquad\text{(D)}\ 3\qquad\text{(E)}\ 6$


 
 \textbf{Answer: }C



\textbf{Solution}

The diagram can be quartered as shown:
                    \begin{figure}[H]
                        centering
                        \includegraphics[width=0.5\textwidth]{Solution/AHSME_I_1989_21_Solution_Figure_1.png}
                    \end{figure}
                    and reassembled into two smaller squares of side $k$, each of which looks like this:
                    \begin{figure}[H]
                        centering
                        \includegraphics[width=0.5\textwidth]{Solution/AHSME_I_1989_21_Solution_Figure_2.png}
                    \end{figure}
                    The border in this figure is the former cross, which still occupies 36\% of the area. Therefore the inner square occupies 64\% of the area, from which we deduce that it is $0.8k \times 0.8k$, and that one blue square must be $0.1k\times 0.1k=0.01k^2$ or 1\% each. Thus the blue area is $\boxed{2\%}$ of the total.

 



\textbf{Solution 2(alcumus)}

Let the side length of the flag be 10, and let each red segment on the border of the flag be $x$. Then the four white triangles can be placed together to form a white square of area $100 - 36 = 64$ and side length $10 - 2x$.

 
                    \begin{figure}[H]
                        centering
                        \includegraphics[width=0.5\textwidth]{Solution/AHSME_I_1989_21_Solution_Figure_3.png}
                    \end{figure}
                    

 Since $(10 - 2x)^2 = 64$, we have $x = 1$. Then the blue area is $(x \sqrt{2})^2 = 2$, which is $\boxed{2\%}$ of the area.

 


	\item (\textit{AHSME 1989}) A child has a set of 96 distinct blocks. Each block is one of 2 materials (plastic, wood), 3 sizes (small, medium, large), 4 colors (blue, green, red, yellow), and 4 shapes (circle, hexagon, square, triangle). How many blocks in the set are different from the 'plastic medium red circle' in exactly 2 ways? (The 'wood medium red square' is such a block)


 (A) 29 (B) 39 (C) 48 (D) 56 (E) 62


 
 \textbf{Answer: }A



\textbf{Solution 1}

The process of choosing a block can be represented by a generating function. Each choice we make can match the 'plastic medium red circle' in one of its qualities $(1)$ or differ from it in $k$ different ways $(kx)$. Choosing the material is represented by the factor $(1+1x)$, choosing the size by the factor $(1+2x)$, etc:
 \[(1+x)(1+2x)(1+3x)^2\]Expanding out the first two factors and the square:
 \[(1+3x+2x^2)(1+6x+9x^2)\]By expanding further we can find the coefficient of $x^2$, which represents the number of blocks differing from the original block in exactly two ways. We don't have to expand it completely, but choose the terms which will be multiplied together to result in a constant multiple of $x^2$:
 \[1\cdot9+3\cdot6+2\cdot1=\boxed{29}\]

 



\textbf{Solution 2}

The blocks can be sorted into two identical $3\times4\times4$ cuboids, one wood and the other plastic, so that in each cuboid the z axis represents the size, the x axis the color, and the y axis the shape.

 Suppose the reference block is in position $(0,0,0)$ in the plastic cuboid.

 The wooden blocks already differ from the reference block in one way, so if $(0,0,0)$ in the wooden cuboid represents 'wood medium red circle' then any wooden block with two zero coordinates satisfies the requirements. These form the edges of the cuboid which adjoin $(0,0,0)$, so there are $2+3+3=8$.

 The required plastic blocks must have two non-zero coordinates, so we count the blocks on the three faces around $(0,0,0)$ but not on the adjoining edges. There are $2\cdot3+2\cdot3+3\cdot3=21$.

 In total there are $29$ blocks meeting the requirements.

 
                    \begin{figure}[H]
                        centering
                        \includegraphics[width=0.5\textwidth]{Solution/AHSME_I_1989_22_Solution_Figure_1.png}
                    \end{figure}
                    

 



\textbf{Solution 3 (Fastest)}

The amount of ways we can differ in the material, sizes, colors, and shapes category is 1, 2, 3, and 3 respectively (we take away one because we cannot choose the characteristic already chosen). We can choose to differ two of these characteristics, so our answer is $1\cdot2 + 1\cdot3 + 1\cdot3 + 2\cdot3 + 2\cdot3 + 3\cdot3 = \boxed{29}$.

 


	\item (\textit{AHSME 1989}) A particle moves through the first quadrant as follows. During the first minute it moves from the origin to $(1,0)$. Thereafter, it continues to follow the directions indicated in the figure, going back and forth between the positive x and y axes, moving one unit of distance parallel to an axis in each minute. At which point will the particle be after exactly 1989 minutes?


 
                    \begin{figure}[H]
                        centering
                        \includegraphics[width=0.5\textwidth]{Problem/AHSME_I_1989_23_Problem_Figure_1.png}
                    \end{figure}
                    


 $\text{(A)}\ (35,44)\qquad\text{(B)}\ (36,45)\qquad\text{(C)}\ (37,45)\qquad\text{(D)}\ (44,35)\qquad\text{(E)}\ (45,36)$


 
 \textbf{Answer: }D



\textbf{Solution}

Squares of size $1\times1,\ 2\times2,\ 3\times3,\ ...$ are successively enclosed between the path and the axes.

 
                    \begin{figure}[H]
                        centering
                        \includegraphics[width=0.5\textwidth]{Solution/AHSME_I_1989_23_Solution_Figure_1.png}
                    \end{figure}
                    

 It takes $1+1+1$ minutes to enclose the first square, $1+2+2$ minutes to enclose the second, $1+3+3$ minutes to enclose the third, and so on. After odd squares, the particle is on the Y axis; after even squares, the particle is on the X axis.

 First we find the highest integer $n$ such that $\sum_{k=1}^n1+2k\le1989$. The sum is equal to
 \[\sum_{k=1}^n1+2\sum_{k=1}^nk=n+n(n+1)=n(n+2)=(n+1)^2-1\]so the highest value is $n=43$ for which the sum is $1935$.

 After $1935$ minutes the 43rd square is enclosed and the particle is at the point $(0,43)$. During the 1936th minute it moves up to $(0,44)$. At the end of the 1980th minute it has moved right to $(44,44)$. After this it moves downward, and at the end of the 1989th minute it is at $(44,35)$. The answer is $\boxed{\textbf{(D)}}$.

 



\textbf{Solution 2}

Note that it takes $n^2+n$ minutes to travel to any point in the format $(n,n)$ <-- aka the top-right corners of the "square-like" shapes. (try this yourself if you dont believe it)

 Additionally, if $n$ is divisible by 2, the next step would be to go down parallel to the y-axis from $(n,n)$. Else, the next step would be to go left parallel to the x-axis from $(n,n)$. We now find a value for $n$ such that $n^2+n = n(n+1) \le 1989$. 

 Note that $40^2 = 1600$. After testing values over 40, we notice that $44(44+1) = 1980$. Therefore it takes 1980 seconds to move to $(44,44)$. Since $n$ is divisible by 2, we just need to move down $9$ units from $(44, 44)$, yielding the final answer of $\boxed{\textbf{(D)}}$.

 


	\item (\textit{AHSME 1989}) Five people are sitting at a round table. Let $f\geq 0$ be the number of people sitting next to at least 1 female and $m\geq0$ be the number of people sitting next to at least one male. The number of possible ordered pairs $(f,m)$ is


 $\mathrm{(A) \ 7 } \qquad \mathrm{(B) \ 8 } \qquad \mathrm{(C) \ 9 } \qquad \mathrm{(D) \ 10 } \qquad \mathrm{(E) \ 11 }$


 
 \textbf{Answer: }B



\textbf{Solution}

Suppose there are more men than women; then there are between zero and two women.

 If there are no women, the pair is $(0,5)$. If there is one woman, the pair is $(2,5)$.

 If there are two women, there are two arrangements: one in which they are together, and one in which they are apart, giving the pairs $(4,5)$ and $(3,4)$.

 All four pairs are asymmetrical; therefore by symmetry there are eight pairs altogether, so $\boxed{\textrm{(B)}}$.

 



\textbf{Solution 2}

Denote $T_n$ as the number of such pairs for $n$ people. Then for $T_{n-1}$, when we add an extra spot, we can either have a male or female giving two options. Note that these two options however double the value of $T_{n-1}$. Now if we note that $T_2=1$, we have that $T_5=8$, so that the answer is $\boxed{\textrm{(B)8}}$.

 


	\item (\textit{AHSME 1989}) In a certain cross country meet between 2 teams of 5 runners each, a runner who finishes in the $n$th position contributes $n$ to his team's score. The team with the lower score wins. If there are no ties among the runners, how many different winning scores are possible?


 (A) 10 (B) 13 (C) 27 (D) 120 (E) 126


 
 \textbf{Answer: }B



\textbf{Solution}

The lowest possible winning score is $1 + 2 + 3 + 4 + 5 = 15$.

 The sum of all $10$ runners' scores is $1 + 2 + \cdots + 10 = 55$.  A team wins if and only if its total score is less than half of this, so the winning team's score is at most $27$. 

 Any score from $15$ to $20$ is possible if the winning team's runners finish in positions $1$, $2$, $3$, $4$, and $x$, where $x$ ranges from $5$ to $10$.

 Any score from $21$ to $25$ is possible if the winning team's runners finish in positions $1$, $2$, $3$, $x$, and $10$, where $x$ ranges from $5$ to $9$.

 Any score from $26$ to $27$ is possible if the winning team's runners finish in positions $1$, $2$, $x$, $9$, and $10$, where $x$ ranges from $4$ to $5$.

 Therefore, any score from $15$ to $27$ is possible, and there are therefore $27 - 15 + 1 = \boxed{\textbf{(B)}\ 13}$ possible winning scores.

 


	\item (\textit{AHSME 1989}) A regular octahedron is formed by joining the centers of adjoining faces of a cube. The ratio of the volume of the octahedron to the volume of the cube is


 $\mathrm{(A) \frac{\sqrt{3}}{12} } \qquad \mathrm{(B) \frac{\sqrt{6}}{16} } \qquad \mathrm{(C) \frac{1}{6} } \qquad \mathrm{(D) \frac{\sqrt{2}}{8} } \qquad \mathrm{(E) \frac{1}{4} }$


 
 \textbf{Answer: }C



\textbf{Solution}

Call the length of a side of the cube x. Thus, the volume of the cube is $x^3$. We can then find that a side of this regular octahedron is the square root of $(\frac{x}{2})^2$+$(\frac{x}{2})^2$ which is equivalent to $\frac{x\sqrt{2}}{2}$. Using our general formula for the volume of a regular octahedron of side length a, which is $\frac{a^3\sqrt2}{3}$, we get that the volume of this octahedron is...

 $(\frac{x\sqrt{2}}{2})^3 \rightarrow \frac{x^3\sqrt{2}}{4} \rightarrow \frac{x^3\sqrt{2}}{4}*\frac{\sqrt{2}}{3} \rightarrow \frac{2x^3}{12}=\frac{x^3}{6}$

 Comparing the ratio of the volume of the octahedron to the cube is\ldots

 $\frac{\frac{x^3}{6}}{x^3} \rightarrow \frac{1}{6}$ or $\fbox{C}$

 



\textbf{Solution 2}

Split the cube in half along a cross-section parallel to two of the bases. Because of symmetry, we only need to compare the remaining square pyramid and the octahedron. Because the area of the base of the pyramid is half the area of the square base of the pyramid, and we divide by three for the pyramid's volume, our final answer is $\frac{1}{2} \cdot \frac{1}{3} = \frac{1}{6}$ or $\fbox{C}$.

 


	\item (\textit{AHSME 1989}) Let $n$ be a positive integer. If the equation $2x+2y+z=n$ has 28 solutions in positive integers $x$, $y$, and $z$, then $n$ must be either 


 $\mathrm{(A)}\ 14 \text{ or } 15\ \mathrm{(B)}\ 15 \text{ or } 16\ \mathrm{(C)}\ 16 \text{ or } 17\ \mathrm{(D)}\ 17 \text{ or } 18\ \mathrm{(E)}\ 18 \text{ or } 19$


 
 \textbf{Answer: }D



\textbf{Solution}

This is equivalent to finding the powers of $k$ with coefficient $28$ in the expansion of $(k^2+k^4+k^6+k^8+...)^2(k+k^2+k^3+k^4+...)$.

 But this is $\left(\frac{k^2}{1-k^2}\right)^2\cdot\frac{k}{1-k}\ =\ k^5\cdot\left(\frac{1}{1-k}\right)^3\cdot\left(\frac{1}{1+k}\right)^2\ =\ k^5\cdot(1+k)\cdot\left(\frac{1}{1-k^2}\right)^3$

 $=k^5(1+k)(\tbinom{2}{2}+\tbinom{3}{2}k^2+\tbinom{4}{2}k^4+\tbinom{5}{2}k^6+...)$, the last part having general term $\tbinom{j}{2}k^{2j-4}$from which it is easy to see that, since $\tbinom{8}{2}=28$, the last part contains the term $28k^{12}$ and the whole result $28k^{17}+28k^{18}$. So the answer is $\mathrm{(D)}$.

 


 



\textbf{Solution 2}

Suppose $n$ is even. Thus $z$ is even, so let $z=2t$. Hence $x+y+t=\frac{n}{2}$. By stars and bars, this has
 \[\binom{\frac{n}{2}-1}{3-1}\]solutions. This implies that $\frac{n}{2}-1=8\to n=18$. If $z$ is odd, let $z=2t-1$. Then $x+y+t=\frac{n+1}{2}$, in which $n=17$ is a solution. The answer is then $\mathrm{(D)}$.

 


	\item (\textit{AHSME 1989}) Find the sum of the roots of $\tan^2x-9\tan x+1=0$ that are between $x=0$ and $x=2\pi$ radians.


 $\mathrm{(A)  \frac{\pi}{2} } \qquad \mathrm{(B) \pi } \qquad \mathrm{(C) \frac{3\pi}{2} } \qquad \mathrm{(D) 3\pi } \qquad \mathrm{(E) 4\pi }$


 
 \textbf{Answer: }D



\textbf{Solution}

The roots of $t^2-9t+1=0$ are positive and distinct, so by considering the graph of $y=\tan x$, the smallest two roots of the original equation $x_1,\ x_2$ are between $0$ and $\frac\pi{2}$, and the two other roots are $\pi+x_1, \pi+x_2$.

 Then, from the quadratic equation, we discover that the product $\tan x_1\tan x_2=1$, which implies that $\tan(x_1+x_2)$ does not exist. The bounds then imply that $x_1+x_2=\frac\pi{2}$. Thus $x_1+x_2+\pi+x_1+\pi+x_2=3\pi$ which is $\textrm{(D)}$.

 



\textbf{Solution 2}

$t^2-9t+1=0$:We treat $\tan(x_1)$ and $\tan(x_2)$ as the roots of our equation. Because $\tan(x_1) \times \tan(x_2) = 1$ by Vieta's formula, $x_1 + x_2 = 0.5\pi$.Because the principal values of $x_1$ and $x_2$ are acute and our range for $x$ is $[0,2\pi]$, we have four values of $x$ that satisfy the quadratic:$x_1, x_2, x_1+\pi, x_2+\pi.$Summing these, we obtain $2(x_1+x_2) + 2\pi$.Using the fact that $x_1+x_2=0.5\pi$,we get $2(0.5\pi) + 2\pi = 3\pi.$

 


	\item (\textit{AHSME 1989}) Find $\sum_{k=0}^{49}(-1)^k\binom{99}{2k}$, where $\binom{n}{j}=\frac{n!}{j!(n-j)!}$


 (A) $-2^{50}$ (B) $-2^{49}$ (C) 0 (D) $2^{49}$ (E) $2^{50}$


 



 
 \textbf{Answer: }B



\textbf{Solution}

By the \href{https://artofproblemsolving.com/wiki/index.php?title=Binomial\_Theorem}{Binomial Theorem}, $(1+i)^{99}=\sum_{n=0}^{99}\binom{99}{j}i^n =$ $\binom{99}{0}i^0+\binom{99}{1}i^1+\binom{99}{2}i^2+\binom{99}{3}i^3+\binom{99}{4}i^4+\cdots +\binom{99}{98}i^{98}$. 

 Using the fact that $i^1=i$, $i^2=-1$, $i^3=-i$, $i^4=1$, and $i^{n+4}=i^n$, the sum becomes: 

 $(1+i)^{99}=\binom{99}{0}+\binom{99}{1}i-\binom{99}{2}-\binom{99}{3}i+\binom{99}{4}+\cdots -\binom{99}{98}$. 

 So, $Re[(1+i)^{99}]=\binom{99}{0}-\binom{99}{2}+\binom{99}{4}-\cdots -\binom{99}{98} = S$. 

 Using \href{https://artofproblemsolving.com/wiki/index.php?title=De\_Moivre\%27s\_Theorem}{De Moivre's Theorem}, $(1+i)^{99}=[\sqrt{2}\mathrm{cis}(45^\circ)]^{99}=\sqrt{2^{99}}\cdot \mathrm{cis}(99\cdot45^\circ)=2^{49}\sqrt{2}\cdot \mathrm{cis}(135^\circ) = -2^{49}+2^{49}i$. 

 And finally, $S=Re[-2^{49}+2^{49}i] = -2^{49}$.

 $\fbox{B}$

 


	\item (\textit{AHSME 1989}) Suppose that 7 boys and 13 girls line up in a row. Let $S$ be the number of places in the row where a boy and a girl are standing next to each other. For example, for the row $GBBGGGBGBGGGBGBGGBGG$ we have that $S=12$. The average value of $S$ (if all possible orders of these 20 people are considered) is closest to


 $\mathrm{(A) 9 } \qquad \mathrm{(B) 10 } \qquad \mathrm{(C) 11 } \qquad \mathrm{(D) 12 } \qquad \mathrm{(E) 13 }$


 
 \textbf{Answer: }A



\textbf{Solution 1}

We approach this problem using Linearity of Expectation. Consider a pair of two people standing next to each other. Ignoring all other people, the probability that a boy is standing on the left position and a girl is standing on the right position is $\frac7{20}\cdot\frac{13}{19}$. Similarly, if a girl is standing on the left position and a boy is standing on the right position the probability is also $\frac{7\cdot 13}{20\cdot 19}$. Thus, the total probability of the two people being one boy and one girl is $\frac{91}{190}$.

 There are a total of 19 different adjacent pairs, so by Linearity of Expectation, we have that the expected value of $S$ is $\frac{91}{10} \to \boxed{A}$.

 



\textbf{Solution 2}

Suppose that the class tried every configuration. Boy $i$ and girl $j$ would stand next to each other in $2$ different orders, in $19$ different positions, $18!$ times each. Summing over all $i,j$ gives $7\cdot13\cdot2\cdot19\cdot18!=\frac{91}{10}\cdot20!$, so the average value of $S$ is $\boxed{9\frac1{10}(A)}$.

 


\end{enumerate}